\documentclass[12pt,letterpaper]{article}

%% ── Packages ─────────────────────────────────────────────────────────────────
\usepackage[margin=1in]{geometry}
\usepackage{amsmath,amssymb,amsthm}
\usepackage{mathtools}
\usepackage{physics}
\usepackage{bm}
\usepackage{graphicx}
\usepackage{booktabs}
\usepackage{array}
\usepackage{tabularx}
\usepackage{multirow}
\usepackage{longtable}
\usepackage{xcolor}
\usepackage{hyperref}
\usepackage{cleveref}
\usepackage{enumitem}
\usepackage{fancyhdr}
\usepackage{titlesec}
\usepackage{titletoc}
\usepackage{tocloft}
\usepackage{caption}
\usepackage{subcaption}
\usepackage{float}
\usepackage{setspace}
%\usepackage{microtype}
\usepackage{cite}
\usepackage{doi}
\usepackage{pgfgantt}
\usepackage[T1]{fontenc}
\usepackage{parskip}
\usepackage{mdframed}
\usepackage{tikz}
\usetikzlibrary{arrows.meta, positioning, shapes.geometric, calc}

%% ── Colour palette ───────────────────────────────────────────────────────────
\definecolor{mitred}{RGB}{163,31,52}
\definecolor{darkgray}{RGB}{60,60,60}
\definecolor{lightgray}{RGB}{245,245,245}
\definecolor{accentblue}{RGB}{0,82,136}
\definecolor{framecolor}{RGB}{163,31,52}

%% ── Hyperref setup ───────────────────────────────────────────────────────────
\hypersetup{
  colorlinks=true,
  linkcolor=mitred,
  citecolor=accentblue,
  urlcolor=accentblue,
  pdftitle={Thesis Proposal: Scaling Laws for Planar Linear Induction Motors},
  pdfauthor={},
}

%% ── Section formatting ───────────────────────────────────────────────────────
\titleformat{\section}{\large\bfseries\color{mitred}}{{\thesection}}{1em}{}[\vspace{-2pt}\rule{\linewidth}{0.4pt}]
\titleformat{\subsection}{\normalsize\bfseries\color{darkgray}}{\thesubsection}{1em}{}
\titleformat{\subsubsection}{\normalsize\itshape\color{darkgray}}{\thesubsubsection}{1em}{}

%% ── Header / footer ──────────────────────────────────────────────────────────
\pagestyle{fancy}
\fancyhf{}
\renewcommand{\headrulewidth}{0.4pt}
\fancyhead[L]{\small\color{darkgray}\textit{Thesis Proposal — Planar LIM Scaling Laws}}
\fancyhead[R]{\small\color{darkgray}EE / AeroAstro Intersection}
\fancyfoot[C]{\small\thepage}

%% ── Theorem environments ─────────────────────────────────────────────────────
\newtheorem{theorem}{Theorem}
\newtheorem{lemma}[theorem]{Lemma}
\newtheorem{proposition}[theorem]{Proposition}
\newtheorem{definition}{Definition}

%% ── Custom environments ──────────────────────────────────────────────────────
\newmdenv[
  linecolor=mitred,
  linewidth=1.5pt,
  backgroundcolor=lightgray,
  leftmargin=0pt,
  rightmargin=0pt,
  innerleftmargin=10pt,
  innerrightmargin=10pt,
  innertopmargin=8pt,
  innerbottommargin=8pt,
]{claimbox}

%% ── Spacing ──────────────────────────────────────────────────────────────────
\setstretch{1.15}
\setlength{\parskip}{6pt}

%% ──────────────────────────────────────────────────────────────────────────────
\begin{document}

%% ══════════════════════════════════════════════════════════════════════════════
%%  TITLE PAGE
%% ══════════════════════════════════════════════════════════════════════════════
\begin{titlepage}
\begin{center}
\vspace*{1cm}

{\color{mitred}\rule{\linewidth}{2pt}}
\vspace{0.4cm}

{\LARGE\bfseries Empirical Derivation of Scaling Laws for\\[6pt]
Planar Printed-Circuit Linear Induction Motors:\\[6pt]
A Miniature Solid-State EMALS Analog}

\vspace{0.4cm}
{\color{mitred}\rule{\linewidth}{2pt}}

\vspace{1.2cm}

{\large\textit{Masters Thesis Proposal}}\\[0.4cm]
{\large Department of Electrical Engineering \& Computer Science}\\[0.2cm]
{\large \& Department of Aeronautics and Astronautics}\\[0.6cm]

\vspace{0.8cm}

\begin{tabular}{rl}
\textbf{Submitted:} & May 17, 2026 \\[4pt]
\textbf{Budget Constraint:} & \$400~USD \\[4pt]
\textbf{Discipline Split:} & 60\% Electrical Engineering ~$|$~ 40\% Aerospace Engineering \\[4pt]
\textbf{Keywords:} & Linear Induction Motor, Magnetic Diffusion, Maxwell Stress \\
                   & Tensor, PCB Electromagnetics, EMALS, Skin Effect, \\
                   & Slip Frequency, Scaling Laws, Finite Element Method \\
\end{tabular}

\vspace{1.5cm}

\begin{mdframed}[linecolor=mitred, linewidth=1pt, backgroundcolor=lightgray]
\centering
\textbf{Central Scientific Claim}\\[6pt]
\textit{This thesis derives and experimentally validates the first closed-form scaling law
relating pole pitch~$\tau_p$, slip frequency~$\omega_s$, and rotor conductivity~$\sigma$
to peak thrust per unit stator area~$\mathcal{F}^*$ for planar PCB Linear Induction Motors
operating in the magnetic skin-effect dominated regime — providing an empirically
grounded design rule directly applicable to sub-100\,W electromagnetic launch systems.}
\end{mdframed}

\vfill

{\small\color{darkgray}
This document constitutes a formal research proposal for a Masters-level investigation
at the intersection of continuum electrodynamics, solid-state power electronics,
and aerospace electromagnetic launch technology.}

\end{center}
\end{titlepage}

%% ══════════════════════════════════════════════════════════════════════════════
%%  TABLE OF CONTENTS
%% ══════════════════════════════════════════════════════════════════════════════
\tableofcontents
\newpage

%% ══════════════════════════════════════════════════════════════════════════════
\section{Introduction and Motivation}
%% ══════════════════════════════════════════════════════════════════════════════

\subsection{Aerospace Context: Electromagnetic Launch Systems}

Electromagnetic launch technology represents one of the most consequential engineering
transitions in modern aerospace. The U.S.\ Navy's Electromagnetic Aircraft Launch System
(EMALS), operational aboard the USS Gerald R.\ Ford (CVN~78), replaced the steam
catapult with a linear induction motor capable of delivering precisely controlled
impulse to aircraft spanning a factor of three in launch weight~\cite{Doyle2010}. Beyond
naval aviation, linear electromagnetic accelerators underpin mass-driver proposals for
lunar surface launch~\cite{Kolm1980}, hyperloop transport, and rapid-fire electromagnetic
ordnance. In all cases, the fundamental physics is identical: a travelling magnetic wave
induces eddy currents in a conductive rotor, and the interaction between induced current
and the applied field produces thrust.

Despite decades of large-scale deployment, the \emph{design-rule literature} for
miniaturised planar geometries operating in the skin-effect dominated regime is
remarkably sparse. Scaling from a 100-tonne aircraft launcher to a 100-gram laboratory
rotor is non-trivial: at reduced scales and elevated spatial frequencies, the magnetic
Reynolds number $\mathrm{Rm} = \mu_0\sigma v L$ and the dimensionless skin depth
$\delta/\tau_p$ enter regimes inaccessible to lumped-circuit LIM models. There are
no published closed-form expressions — validated against physical hardware — that
predict \emph{peak thrust per unit stator area} as a function of the three independently
controllable design parameters: pole pitch, slip frequency, and rotor conductivity.
This thesis fills that gap.

\subsection{Research Gap}

Standard analytical treatments of linear induction motors (Gieras~\cite{Gieras1994};
Boldea \& Nasar~\cite{Boldea2001}) rely on either (a) equivalent-circuit models that
collapse the magnetic skin effect into a single lumped resistance, or (b) two-dimensional
analytical solutions valid only for semi-infinite conductors at low slip. Neither approach
yields a design rule usable by an engineer sizing a planar PCB stator for a compact
electromagnetic launcher. Recent work on PCB motor technology (Lewin et
al.~\cite{Lewin2006}; Raminosoa et al.~\cite{Raminosoa2018}) addresses rotational
geometries and does not treat the travelling-wave, linear-thrust case. The miniature
EMALS analog thus occupies a genuine white space in the literature.

\subsection{Scientific Claim}

\begin{claimbox}
\textbf{Thesis Claim.}\quad There exists a closed-form dimensionless scaling law of the form
\[
  \mathcal{F}^* \;=\; f\!\left(\frac{\delta}{\tau_p},\; \mathrm{Rm}\right)
\]
where $\mathcal{F}^* \equiv F_{\mathrm{thrust}} / (B_0^2 \tau_p w / \mu_0)$ is the
dimensionless thrust, $\delta = \sqrt{2/(\mu_0\sigma\omega_s)}$ is the magnetic skin
depth, and $\mathrm{Rm}$ is the magnetic Reynolds number, whose functional form and
numerical coefficients can be experimentally determined using a parametric sweep over
$\{\tau_p,\,\omega_s,\,\sigma\}$ on a planar PCB LIM platform. This law predicts an
optimal slip frequency $\omega_s^*$ at which thrust is maximised for a given pole pitch,
and that this optimum shifts with rotor conductivity in a manner derivable from the
Convective Magnetic Diffusion equation.
\end{claimbox}

\subsection{Objectives}

The thesis pursues four tightly coupled objectives:

\begin{enumerate}[label=\textbf{O\arabic*.}, leftmargin=2.2cm]
  \item \textbf{Analytical:} Derive the closed-form expression for optimal slip frequency
    $\omega_s^*(\tau_p, \sigma)$ from the Convective Magnetic Diffusion PDE and the
    Maxwell Stress Tensor, including a rigorous proof of the volume-cancellation theorem.
  \item \textbf{Computational:} Develop a validated 2D transient FEM solver (FEniCS) for
    the planar LIM geometry with documented mesh-convergence and parametric sensitivity
    studies across $\tau_p \in [5,\,20]\,$mm and $\omega_s \in [10,\,500]\,$Hz.
  \item \textbf{Experimental:} Fabricate a 4-layer FR4 PCB stator with reconfigurable
    pole pitch and conduct a $3 \times 3 \times 2$ parametric experiment (pole pitch
    $\times$ slip frequency $\times$ rotor alloy) yielding 18 independent data points
    with full uncertainty budgets.
  \item \textbf{Synthesis:} Fit a dimensionless scaling law to the empirical data, compare
    against the analytical prediction, and publish design charts directly applicable to
    sub-100\,W electromagnetic launch systems.
\end{enumerate}

%% ══════════════════════════════════════════════════════════════════════════════
\section{Theoretical Framework}
%% ══════════════════════════════════════════════════════════════════════════════

\subsection{Governing Equation: Convective Magnetic Diffusion PDE}

Starting from Maxwell's equations in the Coulomb gauge $(\nabla \cdot \mathbf{A} = 0)$,
the magnetic flux density is expressed through the magnetic vector potential
$\mathbf{B} = \nabla \times \mathbf{A}$. Inside a moving, isotropic, non-magnetic conductor
of electrical conductivity $\sigma$ and permeability $\mu_0$, travelling at velocity
$\mathbf{v}$ relative to the stator frame, Faraday's law and Ohm's law combine to yield
the Convective Magnetic Diffusion Equation~\cite{Davidson2001}:

\begin{equation}
  \boxed{
  \nabla^2 \mathbf{A} \;=\; \mu_0\sigma\,\frac{\partial \mathbf{A}}{\partial t}
  \;+\; \mu_0\sigma\,(\mathbf{v}\cdot\nabla)\mathbf{A}
  }
  \label{eq:cmd}
\end{equation}

The two terms on the right side of \cref{eq:cmd} encode distinct physics:
the first is \emph{temporal diffusion} (responsible for the magnetic skin effect at
high frequency), and the second is \emph{convective distortion} of the field due to
rotor motion. At low rotor velocity the convective term is negligible and
\cref{eq:cmd} reduces to the standard eddy-current diffusion equation. At high velocity,
the convective term introduces an asymmetric field distribution — compressed on the
leading edge of the rotor, rarefied on the trailing edge — that is the physical origin
of the well-known LIM end-effect efficiency penalty.

\subsection{Dimensionless Formulation and the Magnetic Reynolds Number}

Let $\tau_p$ be the pole pitch (the characteristic length), $\omega_s = \omega - k v_s$
the slip angular frequency (where $\omega$ is the stator excitation frequency and $v_s$
is the synchronous velocity $v_s = \omega \tau_p / \pi$), and $B_0$ the peak air-gap
flux density. Nondimensionalising \cref{eq:cmd} with:
\[
  \tilde{x} = x/\tau_p, \quad
  \tilde{t} = \omega_s t, \quad
  \tilde{\mathbf{A}} = \mathbf{A}/B_0\tau_p,
\]
yields:

\begin{equation}
  \frac{\partial^2 \tilde{A}}{\partial \tilde{x}^2}
  + \frac{\partial^2 \tilde{A}}{\partial \tilde{y}^2}
  \;=\; \underbrace{\mu_0\sigma\omega_s\tau_p^2}_{\displaystyle S}
  \frac{\partial \tilde{A}}{\partial \tilde{t}}
  \;+\; \underbrace{\mu_0\sigma v\tau_p}_{\displaystyle \mathrm{Rm}}\,
  \frac{\partial \tilde{A}}{\partial \tilde{x}}
  \label{eq:cmd_nd}
\end{equation}

where $S = \mu_0\sigma\omega_s\tau_p^2$ is the \emph{quality factor} (dimensionless skin
depth parameter) and $\mathrm{Rm} = \mu_0\sigma v\tau_p$ is the magnetic Reynolds number.
The ratio $S/\mathrm{Rm} = \omega_s\tau_p / v = \pi s$ (where $s$ is the dimensionless slip)
recovers the familiar result from induction machine theory. This nondimensionalisation
reveals that all planar LIM systems with equal $S$ and $\mathrm{Rm}$ are physically
equivalent — the foundation for the proposed scaling law.

\subsection{Magnetic Skin Depth and the Optimal Slip Frequency}

The penetration depth of the travelling field into the rotor is:
\begin{equation}
  \delta \;=\; \sqrt{\frac{2}{\mu_0\sigma\omega_s}}
  \label{eq:skindepth}
\end{equation}

This immediately implies a critical constraint: if the rotor thickness $d < \delta$, eddy
currents penetrate the full cross-section and thrust scales as $\sigma d^2 \omega_s$
(thin-sheet regime). If $d > \delta$, eddy currents are confined to a surface layer and
thrust is independent of $d$ (skin-effect regime). The transition occurs at
$\omega_s^{\mathrm{tr}} = 2/(\mu_0\sigma d^2)$, which for 3\,mm thick 7075-T6 aluminium
($\sigma \approx 1.96 \times 10^7$\,S/m) gives $\omega_s^{\mathrm{tr}} \approx 140$\,Hz.
The proposed experiment spans both regimes deliberately.

\begin{proposition}[Optimal Slip Frequency]
\label{prop:opt_slip}
Under the assumption of a sinusoidal travelling-wave stator field
$B_y = B_0 \cos(\pi x/\tau_p - \omega t)$ and a rotor of half-thickness $d$,
the thrust $F$ as a function of slip frequency attains a maximum at:
\begin{equation}
  \omega_s^* \;=\; \frac{\pi^2}{\mu_0\sigma\tau_p^2}
  \;=\; \frac{\pi^2\,\eta}{\tau_p^2}
  \label{eq:opt_slip}
\end{equation}
where $\eta = 1/(\mu_0\sigma)$ is the magnetic diffusivity. At $\omega_s^*$,
the skin depth satisfies $\delta^* = \sqrt{2}\,\tau_p/\pi$, i.e.\ the skin depth
is proportional to pole pitch, independent of material.
\end{proposition}

\textit{Proof sketch.} Representing the vector potential in the rotor as
$A_z(x,y,t) = \hat{A}(y)\,e^{j(\pi x/\tau_p - \omega_s t)}$ and substituting into
\cref{eq:cmd} (in the rotor rest frame, $\mathbf{v} = 0$), the $y$-dependence satisfies:
\[
  \frac{d^2\hat{A}}{dy^2} \;=\; \left[\left(\frac{\pi}{\tau_p}\right)^2
  - j\mu_0\sigma\omega_s\right]\hat{A} \;\equiv\; \kappa^2\hat{A}
\]
with $\kappa^2 = (\pi/\tau_p)^2(1 - jS)$ where $S = \mu_0\sigma\omega_s\tau_p^2/\pi^2$.
The time-averaged thrust per unit area from the Maxwell Stress Tensor (see \cref{sec:mst})
is proportional to $\mathrm{Im}(\hat{A}^*\partial\hat{A}/\partial y)|_{y=0}$. Maximising
over $S$ yields $S^* = 1$, i.e.\ $\omega_s^* = \pi^2/(\mu_0\sigma\tau_p^2)$. \hfill$\square$

\Cref{prop:opt_slip} is the central falsifiable prediction of this thesis. It predicts that
a doubling of pole pitch reduces the optimal slip frequency by a factor of four,
and that changing from 7075-T6 aluminium to 6061-T6 (a 15\% reduction in conductivity)
shifts $\omega_s^*$ upward by 15\%. Both predictions will be tested experimentally.

\subsection{Volume Cancellation Theorem}
\label{sec:vol_cancel}

A preliminary result from the 1D kinematic model warrants formal statement because it
has a direct bearing on experimental design.

\begin{theorem}[Volume Cancellation]
\label{thm:vol_cancel}
Under idealized uniform magnetic diffusion (thin-sheet limit, $d \ll \delta$),
the equation of motion for the rotor reduces to:
\begin{equation}
  \dot{v} \;=\; \frac{\sigma B_0^2 \omega_s d}{\rho}\,g(s)
  \label{eq:vol_cancel}
\end{equation}
where $\rho$ is the mass density, $d$ the rotor thickness, and $g(s)$ a dimensionless
function of slip. In particular, the rotor volume $V = A_{\mathrm{contact}} \cdot d$
cancels: acceleration is independent of rotor plan-area and depends on thickness $d$
only through the ratio $\sigma d / \rho$.
\end{theorem}

\begin{proof}
The induced EMF per unit length in the thin-sheet rotor is $\mathcal{E} = v_s B_0 s$
(slip times synchronous velocity). The short-circuit current density is
$J = \sigma \mathcal{E} / d$ (resistance of a slab of thickness $d$). The Lorentz
force per unit volume is $f = J B_0 = \sigma v_s s B_0^2 / d$. Newton's second law
for the rotor of mass $m = \rho A d$:
\[
  \rho A d\, \dot{v} \;=\; f \cdot A d \;=\; \sigma v_s s B_0^2 A
\]
The $A$ and $d$ cancel, giving $\dot{v} = \sigma v_s s B_0^2 / \rho$,
confirming \cref{eq:vol_cancel}. Note $\rho$ appears but volume does not. \end{proof}

\textbf{Experimental implication.} \Cref{thm:vol_cancel} implies that in the thin-sheet
regime, rotor mass need not be controlled between experimental runs — only the
$\sigma/\rho$ ratio matters. For 7075-T6 and 6061-T6, $\sigma/\rho$ differs by
approximately 10\%, providing a meaningful experimental contrast with nearly identical
mechanical behaviour.

\subsection{Force Calculation: Maxwell Stress Tensor}
\label{sec:mst}

Rather than computing thrust from idealized uniform eddy current distributions,
this project integrates the Maxwell Stress Tensor~\cite{Griffiths2017} over a closed
surface $\partial\mathcal{V}$ enclosing the rotor:

\begin{equation}
  \overleftrightarrow{T}_{ij} \;=\; \frac{1}{\mu_0}
  \left[B_i B_j \;-\; \frac{1}{2}\,\delta_{ij}\,B^2\right]
  \label{eq:mst}
\end{equation}

The time-averaged propulsive thrust per unit stator width $w$ is:

\begin{equation}
  \langle F_x \rangle \;=\; w \oint_{\partial\mathcal{V}}
  \langle T_{xj}\rangle\, n_j\, dS
  \;=\; w \int_0^{2\tau_p}
  \frac{\langle B_x B_y\rangle}{\mu_0}\Bigg|_{y=0}\!\! dx
  \label{eq:thrust}
\end{equation}

where $y = 0$ is the stator surface and the surface integral reduces to the
air-gap tangential stress. For the analytical sinusoidal solution:

\begin{equation}
  \langle F_x \rangle \;=\; \frac{B_0^2 \tau_p w}{2\mu_0}\,
  \frac{S}{1 + S^2}
  \label{eq:thrust_analytical}
\end{equation}

\Cref{eq:thrust_analytical} — which peaks at $S = 1$, recovering \cref{eq:opt_slip} —
is the primary theoretical prediction to be validated by FEM and experiment.
The dimensionless thrust is thus:

\begin{equation}
  \boxed{\mathcal{F}^* \;=\; \frac{S}{1 + S^2}, \qquad
  S \;=\; \frac{\mu_0\sigma\omega_s\tau_p^2}{\pi^2}}
  \label{eq:scaling_law}
\end{equation}

\Cref{eq:scaling_law} is the closed-form scaling law this thesis will validate.

%% ══════════════════════════════════════════════════════════════════════════════
\section{Literature Review and Position}
%% ══════════════════════════════════════════════════════════════════════════════

\subsection{Classical LIM Theory}

Gieras~\cite{Gieras1994} established the foundational analytical treatment of linear
induction machines, deriving thrust-slip curves from a 2D magnetoquasistatic model.
Boldea and Nasar~\cite{Boldea2001} extended this to include end effects and dynamic braking.
Both frameworks use lumped equivalent circuits with frequency-dependent parameters,
which capture skin-effect trends but do not yield closed-form design rules for the
$\tau_p$--$\omega_s$--$\sigma$ parameter space.

\subsection{EMALS and Large-Scale Linear Actuators}

The Navy's EMALS program~\cite{Doyle2010} represents the state of the art in linear
electromagnetic launch. Connor et al.~\cite{Connor2015} document the power electronics
and energy storage architecture. Neither source addresses miniaturisation or PCB-scale
realisation. Kolm and Mongeau~\cite{Kolm1980} analyzed electromagnetic mass drivers for
space launch but operated in a fundamentally different regime (pulsed, non-inductive
thrust) from the continuous-wave LIM addressed here.

\subsection{PCB Motor Technology}

Lewin et al.~\cite{Lewin2006} demonstrated that multi-layer PCB windings can achieve
sufficient current density for low-power actuators. Raminosoa et al.~\cite{Raminosoa2018}
characterised thermal limits of PCB stator coils under sustained operation. Neither work
addresses the linear (translational) geometry or the skin-effect dominated thrust regime.
Maloberti et al.~\cite{Maloberti2020} simulated planar eddy-current brakes using FEM
but did not investigate the thrust-maximising slip frequency experimentally.

\subsection{Identified Gap}

No published work provides: (a) an analytically derived, (b) computationally validated,
and (c) physically confirmed closed-form scaling law for peak thrust per unit area as a
function of $\{\tau_p,\,\omega_s,\,\sigma\}$ in a planar PCB LIM operating in the
skin-effect dominated regime. This thesis provides exactly that.

%% ══════════════════════════════════════════════════════════════════════════════
\section{Numerical Simulation Methodology}
%% ══════════════════════════════════════════════════════════════════════════════

The simulation strategy is deliberately three-tiered, with each tier trading
fidelity for speed. The 1D ODE solver (complete) drives parameter-space exploration;
the 2D FEM on Google Colab provides rapid validation of the analytical scaling law
at low cost; and the 3D FEM on a local Mac Mini M4 captures end effects and flux
fringing that the 2D model structurally cannot represent. The two FEM tiers are
complementary, not redundant.

\subsection{Compute Platform Summary}

\begin{table}[H]
\centering
\caption{Simulation compute platforms and realistic constraints.}
\label{tab:compute}
\begin{tabularx}{\textwidth}{lXll}
\toprule
\textbf{Platform} & \textbf{Specs} & \textbf{Use} & \textbf{Cost} \\
\midrule
Python / laptop & Any CPU, \texttt{scipy} & Tier 1 ODE & Free \\
Google Colab T4  & 15\,GB VRAM, $\sim$5\,hr/day free (unguaranteed) & Tier 2 2D FEM & Free \\
Mac Mini M4      & Apple M4, 16\,GB unified RAM, macOS & Tier 3 3D FEM & \$0 (owned) \\
\bottomrule
\end{tabularx}
\end{table}

\noindent\textbf{Note on Colab constraints.} Google Colab free tier provides T4 GPU
access for approximately 5 hours per day, subject to availability and session
disconnection. All 2D FEM jobs are therefore designed to complete in under 90 minutes
per run, with checkpointing to \texttt{.h5} files so that a disconnected session can
resume mid-sweep. No A100 or paid compute is assumed anywhere in this proposal.

\subsection{Tier 1 — 1D Kinematic ODE Solver (Completed)}

A Python solver using \texttt{scipy.integrate.solve\_ivp} integrates the rotor
equation of motion derived in \cref{thm:vol_cancel}. The solver accepts
$\{B_0,\,\tau_p,\,\omega_s,\,\sigma,\,\rho,\,d\}$ as inputs and outputs position
$x(t)$, velocity $v(t)$, and instantaneous thrust $F(t)$. It has been validated
against the analytical impulse prediction for the thin-sheet limit and reproduces the
volume-cancellation result numerically. It serves as the fast surrogate for
parameter-space exploration, guiding which operating points are most informative
before committing FEM resources.

The governing ODE system is:
\begin{align}
  \dot{x} &= v \label{eq:ode1}\\
  \dot{v} &= \frac{F_{\mathrm{thrust}}(v,\,\omega_s,\,\tau_p,\,\sigma)}{m}
             - \frac{F_{\mathrm{drag}}(v)}{m} \label{eq:ode2}
\end{align}

where $F_{\mathrm{thrust}}$ is evaluated from \cref{eq:thrust_analytical} with slip
$s = 1 - v\pi/(\omega\tau_p)$ updated at each timestep, and $F_{\mathrm{drag}}$
accounts for aerodynamic and rail-contact drag.

\textbf{Estimated runtime:} $< 2$\,seconds per trajectory on any modern laptop.
Full 18-point parametric sweep: $< 1$\,minute.

\subsection{Tier 2 — 2D Transient FEM on Google Colab T4}
\label{sec:tier2}

\subsubsection{Motivation and Scope}

The 2D transient FEM solves \cref{eq:cmd} in the cross-sectional plane
$(x,\,y)$ — the direction of rotor travel and the depth into the rotor —
capturing the magnetic skin effect and its dependence on $\omega_s$, $\tau_p$,
and $\sigma$ with high fidelity. It validates the analytical scaling law
\cref{eq:scaling_law} under realistic boundary conditions and is the primary
tool for the mesh convergence study. Its fundamental limitation is that it
assumes the rotor is infinitely wide in the $z$-direction (transverse to motion),
eliminating edge effects. This limitation is addressed by Tier~3.

\subsubsection{Implementation}

The solver uses FEniCS~\cite{AlnaesEtal2015} (Python UFL formulation, legacy
2019.1 version, pip-installable on Colab without sudo). The computational domain
is a two-region rectangle:

\begin{itemize}[noitemsep]
  \item \textbf{Air-gap + stator region} ($0 \leq y \leq g$): prescribed sinusoidal
    surface current $K_z = K_0\cos(\pi x/\tau_p - \omega t)$ at $y = 0$.
  \item \textbf{Rotor region} ($-d \leq y \lewidth 0$): \cref{eq:cmd} governs;
    material parameters set to 7075-T6 or 6061-T6 aluminium.
\end{itemize}

Boundary conditions: $A_z = 0$ on top/bottom (magnetic insulation, domain height
$= 10\delta$); periodic left/right with phase offset $\Delta\phi = 2\pi\tau_p/\lambda$.
Time integration uses a second-order Crank--Nicolson scheme with
$\Delta t = 1/(20\,\omega_s)$ (20 steps per slip cycle), run for 5 slip cycles to
reach periodic steady state before extracting thrust.

\subsubsection{Mesh Convergence Protocol}

Three structured quadrilateral meshes will be generated in the rotor region:

\begin{table}[H]
\centering
\caption{2D FEM mesh convergence study parameters and estimated runtimes (Colab T4).}
\label{tab:mesh2d}
\begin{tabular}{lccccc}
\toprule
\textbf{Mesh} & \textbf{Element size} & \textbf{DOF (approx.)} & \textbf{Time steps} & \textbf{Runtime / point} & \textbf{Total (30 pts)} \\
\midrule
Coarse & $h = \delta/2$  & $\sim$8\,k   & 100 & $\sim$4\,min  & $\sim$2\,hr \\
Medium & $h = \delta/4$  & $\sim$30\,k  & 100 & $\sim$14\,min & $\sim$7\,hr \\
Fine   & $h = \delta/8$  & $\sim$110\,k & 100 & $\sim$50\,min & $\sim$25\,hr \\
\bottomrule
\end{tabular}
\end{table}

\noindent The convergence study uses the medium mesh for the full parametric sweep
(30 points $\times$ 14\,min $\approx$ 7 hours total, spread across $\sim$2 Colab
sessions). The fine mesh is run only at the 3 operating points closest to $S = 1$
to confirm Richardson extrapolation bounds. All jobs checkpoint to Google Drive after
each operating point so session disconnection causes at most one lost point.

Convergence is declared when thrust changes by $< 2\%$ between successive refinements.

\subsubsection{Parametric Sweep}

The 2D FEM evaluates \cref{eq:scaling_law} across
$\tau_p \in \{5,\,10,\,20\}$\,mm and $\omega_s/2\pi \in \{10,\,50,\,100,\,200,\,500\}$\,Hz
for both rotor alloys (30 operating points total). Results are post-processed to
extract $\mathcal{F}^*$ and $S$ and overlaid on the analytical curve.

\subsection{Tier 3 — 3D Transient FEM on Mac Mini M4}
\label{sec:tier3}

\subsubsection{Why 3D is Necessary}

The 2D model assumes a rotor of infinite transverse width ($w \to \infty$). In the
physical experiment the rotor width is 12\,mm — comparable to the pole pitch at the
5\,mm and 10\,mm configurations. Eddy currents that the 2D model assumes flow purely
in the $x$--$y$ plane must in reality close through the $z$-direction, constrained by
the rotor's finite width. This introduces two effects absent from the 2D model:

\begin{enumerate}[noitemsep]
  \item \textbf{Transverse edge flux leakage}: the field bulges outward at the rotor
    edges, reducing the effective flux linkage and therefore thrust. This penalty is
    largest when $\tau_p \sim w$ (the 5\,mm pole-pitch configuration).
  \item \textbf{3D eddy current path elongation}: eddy current loops must travel
    around the rotor perimeter in $z$, increasing their effective resistance and
    reducing current magnitude for a given induced EMF.
\end{enumerate}

Both effects introduce a geometry-dependent correction factor $\eta_{3D}(\tau_p/w)$
that multiplies the 2D thrust prediction. Quantifying $\eta_{3D}$ is a genuine
scientific contribution beyond the 2D result.

\subsubsection{Software Stack: Elmer FEM + Gmsh}

The 3D solver uses \textbf{Elmer FEM}~\cite{ElmerTeam2020} — an open-source,
actively maintained finite element package from CSC Finland with a dedicated
\texttt{MagnetoDynamics} solver module implementing the $\mathbf{A}$--$V$ formulation
of \cref{eq:cmd} in full 3D. Elmer runs natively on macOS ARM (Apple Silicon) via
Homebrew and has been benchmarked on M-series hardware by the community.
Mesh generation uses \textbf{Gmsh}~\cite{Geuzaine2009} (also Homebrew-installable),
which produces structured hexahedral meshes in the rotor bulk and unstructured
tetrahedral meshes in the air-gap and fringe regions.

The full toolchain is free and open-source. No licences, no cloud costs.

\subsubsection{3D Domain and Mesh}

The 3D computational domain replicates the physical rotor geometry:
$150\,\mathrm{mm} \times 12\,\mathrm{mm} \times d\,\mathrm{mm}$ rotor (where
$d \in \{1,\,2,\,3\}$\,mm), suspended above a 1\,mm air gap over the stator surface.
The stator is represented as a prescribed current-sheet boundary condition
(surface current $\mathbf{K}$), avoiding the need to mesh the PCB copper traces
in 3D. The domain is truncated at $5\delta$ above and below with homogeneous
Dirichlet conditions $\mathbf{A} = \mathbf{0}$.

\begin{table}[H]
\centering
\caption{3D FEM mesh parameters and estimated runtimes on Mac Mini M4 (16\,GB unified RAM).}
\label{tab:mesh3d}
\begin{tabular}{lcccc}
\toprule
\textbf{Config.} & \textbf{DOF (approx.)} & \textbf{RAM usage} & \textbf{Runtime / point} & \textbf{Parallelism} \\
\midrule
$\tau_p = 20$\,mm, $d = 3$\,mm (coarsest) & $\sim$400\,k & $\sim$4\,GB & $\sim$2\,hr & 4 threads (P-cores) \\
$\tau_p = 10$\,mm, $d = 2$\,mm (medium)   & $\sim$900\,k & $\sim$9\,GB & $\sim$4\,hr & 4 threads \\
$\tau_p = 5$\,mm, $d = 1$\,mm (finest)    & $\sim$1.8\,M & $\sim$15\,GB & $\sim$8\,hr & 4 threads \\
\bottomrule
\end{tabular}
\end{table}

\noindent\textbf{Runtime strategy.} The three configurations in \cref{tab:mesh3d}
represent the 3 pole-pitch levels at a single operating frequency
($\omega_s = \omega_s^*$, the analytically predicted optimum). Each run executes
overnight on the Mac Mini M4, which has no duty-cycle or session-time constraints.
Total 3D compute time for the primary result: approximately 14 hours of wall time
across 3 overnight runs. An additional 6 configurations (2 alloys $\times$ 3 pitches
at a second frequency point) add $\sim$30 hours, comfortably within a 2-week
simulation window.

The M4's 16\,GB unified memory architecture is critical: at 1.8M DOFs the solver
requires $\sim$15\,GB, which fits in unified RAM but would overflow a discrete GPU's
VRAM on a cheaper machine. Elmer's OpenMP parallelism uses all 4 performance cores
of the M4; efficiency-cores are not used (Elmer does not currently support
heterogeneous thread scheduling on Apple Silicon).

\subsubsection{3D Parametric Scope}

The 3D sweep is deliberately narrower than the 2D sweep, targeting the correction
factor $\eta_{3D}$ rather than re-tracing the full $\mathcal{F}^*$ vs.\ $S$ curve:

\begin{itemize}[noitemsep]
  \item \textbf{Primary:} 3 pole pitches $\times$ 2 alloys $\times$ 2 frequency
    points = 12 operating points. \textit{Estimated total runtime: 36\,hours (6 overnight sessions).}
  \item \textbf{Subsidiary:} 3 rotor thicknesses at fixed $\tau_p = 10$\,mm for
    Volume Cancellation Theorem validation (3 points $\times$ 4\,hr = 12\,hours).
\end{itemize}

\textbf{Total 3D FEM wall time: $\sim$48\,hours across $\sim$8 overnight runs on the
Mac Mini M4.} No single run exceeds 8 hours, making overnight scheduling reliable.

\subsubsection{Deliverable from 3D FEM}

The primary 3D output is the empirical correction factor:
\begin{equation}
  \eta_{3D}(\tau_p/w) \;=\;
  \frac{\langle F_x\rangle_{\mathrm{3D}}}{\langle F_x\rangle_{\mathrm{2D}}}
  \Bigg|_{\omega_s = \omega_s^*}
  \label{eq:eta3d}
\end{equation}
plotted as a function of the aspect ratio $\tau_p/w$. If $\eta_{3D}$ is found to
follow a simple functional form (e.g.\ $1 - c(\tau_p/w)^n$), this constitutes an
additional closed-form result beyond \cref{eq:scaling_law} — a transverse-edge
correction directly usable in PCB LIM design.

\subsection{Validation of Volume Cancellation Theorem by FEM}

A subsidiary 2D FEM study (Tier~2, $\sim$45\,min total runtime on Colab T4)
verifies \cref{thm:vol_cancel} by simulating three rotor thicknesses
($d \in \{1,\,2,\,3\}$\,mm) at fixed $\omega_s < \omega_s^{\mathrm{tr}}$ and
confirming that computed thrust scales as $F \propto d$ (thin-sheet limit), not
$d^2$ or $d^0$. The 3D FEM runs the same study to check whether the $z$-edge
correction modifies this scaling in the finite-width geometry.

%% ══════════════════════════════════════════════════════════════════════════════
\section{Experimental Design}
%% ══════════════════════════════════════════════════════════════════════════════

\subsection{Design Philosophy: Parametric Coverage Over Single-Point Demonstration}

The defining upgrade from demonstration to thesis is the shift from one operating
point to a \emph{design space map}. The experiment is structured as a
$3 \times 3 \times 2$ factorial design:

\begin{table}[H]
\centering
\caption{Experimental factorial design (18 independent runs).}
\label{tab:factorial}
\begin{tabular}{llll}
\toprule
\textbf{Factor} & \textbf{Level 1} & \textbf{Level 2} & \textbf{Level 3} \\
\midrule
Pole pitch $\tau_p$ & 5\,mm & 10\,mm & 20\,mm \\
Slip frequency $\omega_s/2\pi$ & 50\,Hz & 150\,Hz & 400\,Hz \\
Rotor alloy & 7075-T6 Al & 6061-T6 Al & — \\
\bottomrule
\end{tabular}
\end{table}

The three pole-pitch levels span a factor of four in $S$ (at fixed $\omega_s$),
crossing the peak of the thrust--slip curve twice. The three slip frequencies straddle
$\omega_s^*$ for 10\,mm pitch / 7075 alloy. The two rotor alloys differ in conductivity
by $\sim$15\%, providing a controlled material contrast. Together the 18 runs will
trace the full curve of \cref{eq:scaling_law} in dimensionless coordinates.

\subsection{PCB Stator Architecture}

\subsubsection{Layer Stack-up}

The stator is a custom 4-layer FR4 PCB (1.6\,mm total, standard JLCPCB 4-layer
stackup: \$35 per board including shipping). Layer allocation:

\begin{table}[H]
\centering
\caption{4-layer PCB stackup and function.}
\label{tab:pcb}
\begin{tabular}{cll}
\toprule
\textbf{Layer} & \textbf{Copper weight} & \textbf{Function} \\
\midrule
L1 (top)    & 2\,oz & Phase~A meandered coil traces \\
L2          & 1\,oz & Phase~B meandered coil traces \\
L3          & 1\,oz & Phase~C meandered coil traces \\
L4 (bottom) & 1\,oz & Control electronics and power distribution \\
\bottomrule
\end{tabular}
\end{table}

\subsubsection{Pole Pitch Configurability}

To access three pole-pitch values within a \$400 budget, a \emph{jumper-configurable
pole connection} scheme will be implemented: a single PCB layout contains trace
patterns for all three pitches, with solder jumpers selecting which traces are
electrically active. This eliminates the cost of ordering three separate PCBs
($\sim$\$35 vs.\ $\sim$\$105).

\subsubsection{Trace Design from PDE Constraints}

The spatial wavenumber of the excited travelling wave must match the pole pitch:
$k = \pi/\tau_p$. For a sinusoidal current distribution, the optimal trace width
$w_{\mathrm{trace}}$ and spacing $g$ satisfy:
\begin{equation}
  w_{\mathrm{trace}} + g \;=\; \frac{\tau_p}{N_{\mathrm{turns}}}
  \label{eq:trace_design}
\end{equation}
where $N_{\mathrm{turns}}$ is the number of turns per pole. The minimum trace width is
set by the current-carrying capacity of 2\,oz copper: at $I_{\mathrm{peak}} = 3$\,A,
a minimum trace width of 0.8\,mm is required (IPC-2221, $\Delta T = 10$\,°C).
The final trace layout will be designed in KiCad~8 and exported as Gerber files for
JLCPCB PCBA manufacturing.

\subsection{Power Electronics and Microcontroller}

\subsubsection{3-Phase Sinusoidal PWM Generation}

The RP2040 microcontroller generates three-phase SPWM using its Programmable I/O
(PIO) state machines operating at 125\,MHz. Each PIO executes a lookup-table-based
sine waveform with 256 steps per period, achieving better than 0.5\% total harmonic
distortion (THD) at frequencies up to 1\,kHz. The three channels are phase-shifted by
$\Delta\phi = 2\pi/3$ via index offsets in the shared lookup table.

The slip frequency is set by software: given stator frequency $\omega$ and rotor
velocity $v$ (updated from DAQ), the RP2040 computes $\omega_s = \omega - \pi v/\tau_p$
and adjusts the PWM period in real time. This closed-loop slip control is a key
capability not present in the original proposal.

\subsubsection{Gate Driver and Power Stage}

The DRV8313 3-phase gate driver bridges the RP2040 logic levels to the 24\,V stator
supply. Maximum continuous output current is 2.5\,A per phase (peak 5\,A for
$< 1$\,s). A bootstrap circuit on each high-side switch eliminates the need for an
isolated gate supply. Dead-time insertion is set to 500\,ns to prevent shoot-through.
Current sensing is implemented via 0.1\,$\Omega$ shunt resistors with the INA240A
current sense amplifier, providing real-time phase current measurement for power
accounting and efficiency calculation.

\subsection{Kinematic Data Acquisition System}

\subsubsection{Multi-Modal Measurement}

A single measurement modality (photointerrupter) carries all kinematic information in
the original proposal. The thesis adds a second independent modality to enable
uncertainty cross-validation:

\begin{itemize}[noitemsep]
  \item \textbf{Optical photointerrupters} (TCST2103): 8 sensors at 15\,mm spacing
    along the launch rail. RP2040 timestamps beam breaks with 8\,ns resolution
    (125\,MHz system clock).
  \item \textbf{Reflective optical encoder strip}: A 500 lines/inch adhesive encoder
    strip on the rotor edge, read by an AS5311 linear magnetic encoder via SPI.
    Provides continuous position at 1\,$\mu$m resolution, enabling direct velocity
    differentiation rather than finite-difference between sensors.
\end{itemize}

Both modalities output position-versus-time traces. Their agreement (or measured
discrepancy) provides an experimental cross-check independent of any model assumption.

\subsubsection{Uncertainty Budget}

\begin{table}[H]
\centering
\caption{Measurement uncertainty budget for kinematic DAQ.}
\label{tab:uncertainty}
\begin{tabularx}{\textwidth}{Xcccc}
\toprule
\textbf{Source} & \textbf{Type} & \textbf{Magnitude} & \textbf{Sensitivity} & \textbf{$u(F)$ contribution} \\
\midrule
Timestamp resolution & B & 8\,ns & $\partial v/\partial t$ & $<0.01$\,mm/s \\
Sensor position tolerance & B & $\pm 0.1$\,mm & $\partial v/\partial x$ & $0.05$\% \\
Rotor temperature ($\sigma$ drift) & B & $\pm 2$\,°C & $\partial\sigma/\partial T$ & $0.3$\% \\
Stray magnetic coupling & A & estimated & Monte Carlo & $< 1$\% \\
Trigger jitter (TCST2103) & A & $\pm 2\,\mu$s & $\partial v/\partial t$ & $0.02$\,mm/s \\
\midrule
\textbf{Combined ($k=2$)} & & & & $< 2$\% \\
\bottomrule
\end{tabularx}
\end{table}

The combined uncertainty of $< 2\%$ on velocity (and thus $< 4\%$ on inferred
kinetic energy / thrust) is sufficient to discriminate the $15\%$ conductivity
difference between the two rotor alloys and to resolve the $\sim 30\%$ thrust
reduction predicted at $\omega_s = 2\omega_s^*$ vs.\ $\omega_s^*$.

\subsection{Thrust Inference from Kinematics}

Direct force measurement requires a load cell, which is incompatible with the free-flight
rotor geometry. Instead, thrust is inferred from kinematics using two methods:

\begin{enumerate}[noitemsep]
  \item \textbf{Instantaneous power method}: $F(t) = P_{\mathrm{mech}}(t) / v(t)$,
    where $P_{\mathrm{mech}}$ is the mechanical power computed from the time derivative
    of kinetic energy: $P_{\mathrm{mech}} = mv\dot{v}$.
  \item \textbf{Impulse method}: $\langle F \rangle = \Delta(mv) / \Delta t$, averaged
    over the acceleration interval. This is model-independent and requires only
    initial and final velocity.
\end{enumerate}

Consistency between the two methods, and between the two DAQ modalities, provides a
four-way experimental cross-check without any additional hardware cost.

%% ══════════════════════════════════════════════════════════════════════════════
\section{Budget and Bill of Materials}
%% ══════════════════════════════════════════════════════════════════════════════

\begin{table}[H]
\centering
\caption{Complete bill of materials (\$400 budget).}
\label{tab:bom}
\begin{tabularx}{\textwidth}{Xlrr}
\toprule
\textbf{Item} & \textbf{Supplier} & \textbf{Qty} & \textbf{Cost (USD)} \\
\midrule
\multicolumn{4}{l}{\textit{PCB and Manufacturing}} \\
4-layer FR4 PCB + PCBA (JLCPCB) & JLCPCB & 5 boards & \$55 \\
\midrule
\multicolumn{4}{l}{\textit{Power Electronics}} \\
DRV8313 3-Phase Gate Driver & DigiKey & 3 & \$12 \\
INA240A Current Sense Amplifier & DigiKey & 6 & \$14 \\
24V 5A DC Power Supply (bench) & Amazon & 1 & \$28 \\
0.1$\Omega$ shunt resistors (2W) & DigiKey & 20 & \$6 \\
Bootstrap capacitors, gate resistors & DigiKey & kit & \$8 \\
\midrule
\multicolumn{4}{l}{\textit{Microcontroller and Logic}} \\
Raspberry Pi Pico (RP2040) & Adafruit & 3 & \$15 \\
AS5311 linear magnetic encoder IC & DigiKey & 2 & \$18 \\
TCST2103 photointerrupters & DigiKey & 12 & \$14 \\
\midrule
\multicolumn{4}{l}{\textit{Rotor Materials}} \\
7075-T6 aluminium bar (12$\times$3$\times$150\,mm) & McMaster & 4 pieces & \$22 \\
6061-T6 aluminium bar (12$\times$3$\times$150\,mm) & McMaster & 4 pieces & \$18 \\
500 lines/inch encoder strip (adhesive) & US Digital & 0.5\,m & \$24 \\
\midrule
\multicolumn{4}{l}{\textit{Mechanical and Fixturing}} \\
3D printed launch rail (PETG) & Local/Bambu & 1 set & \$12 \\
M3 hardware, standoffs, connectors & Amazon & kit & \$14 \\
Breadboard + jumper wires & generic & 2 & \$10 \\
\midrule
\multicolumn{4}{l}{\textit{Test and Measurement}} \\
Logic analyser (16-ch, USB) & Saleae clone & 1 & \$22 \\
Thermocouple + reader (rotor temp) & Amazon & 1 & \$16 \\
\midrule
\multicolumn{4}{l}{\textit{Contingency}} \\
Component replacement / shipping & — & — & \$32 \\
\midrule
\textbf{Total} & & & \textbf{\$340} \\
\textbf{Budget remaining (margin)} & & & \textbf{\$60} \\
\bottomrule
\end{tabularx}
\end{table}

\textbf{Note on compute costs.} All FEM simulation runs on Google Colab free tier
(T4 GPU, 15\,GB VRAM). For the parametric sweep, estimated runtime is
$\sim$4 hours total — well within Colab's daily free allocation.

%% ══════════════════════════════════════════════════════════════════════════════
\section{Phase-Wise Project Roadmap}
%% ══════════════════════════════════════════════════════════════════════════════

\subsection{Overview}

The project is structured in five sequential phases. Phases 1 is complete.
Phases 2 through 5 constitute the outstanding work.

\begin{figure}[H]
\centering
\begin{ganttchart}[
  hgrid,
  vgrid={*{3}{dotted}, *1{dashed}},
  title label font=\small,
  bar label font=\small,
  milestone label font=\small\bfseries,
  bar/.style={fill=mitred!70},
  bar incomplete/.style={fill=mitred!25},
  milestone/.style={fill=mitred, draw=mitred},
  bar height=0.55,
  y unit chart=0.65cm,
  x unit=0.55cm,
  time slot format=isodate-yearmonth
]{2026-05}{2026-12}
  \gantttitlecalendar{year, month=shortname} \\
  \ganttbar[bar/.style={fill=accentblue!60}]{\textbf{Ph.1} Theory (Done)}{2026-05}{2026-05} \\
  \ganttbar{\textbf{Ph.2} PDE / FEM}{2026-06}{2026-07} \\
  \ganttmilestone{FEM validated}{2026-07} \\
  \ganttbar{\textbf{Ph.3} Hardware Design}{2026-07}{2026-08} \\
  \ganttmilestone{Gerbers dispatched}{2026-08} \\
  \ganttbar{\textbf{Ph.4} Firmware + Integration}{2026-08}{2026-09} \\
  \ganttmilestone{First rotor launch}{2026-09} \\
  \ganttbar{\textbf{Ph.5} Parametric Experiment}{2026-10}{2026-11} \\
  \ganttmilestone{Thesis draft}{2026-12} \\
\end{ganttchart}
\caption{Project Gantt chart (May -- December 2026).}
\label{fig:gantt}
\end{figure}

\subsection{Phase 1 — Theoretical Foundation (Complete \checkmark)}

\textbf{Duration:} Completed prior to submission.

\textbf{Deliverables achieved:}
\begin{itemize}[noitemsep]
  \item Idealized Lorentz force equations established; F = J$\times$B framework derived.
  \item Volume Cancellation Theorem proved analytically (see \cref{thm:vol_cancel}).
  \item GitHub repository initialised: Python 1D kinematic ODE solver operational.
  \item Three-phase waveform visualizer validated against analytical SPWM spectra.
  \item Preliminary parameter-space survey identifying the $S$-parameter as the
    governing dimensionless group.
\end{itemize}

\textbf{Key finding:} The dimensionless quality factor $S = \mu_0\sigma\omega_s\tau_p^2/\pi^2$
fully parameterises the analytical thrust curve (\cref{eq:scaling_law}), collapsing all
material and geometric variation into a single curve. This motivates the experimental
design strategy.

\subsection{Phase 2 — Rigorous Analytical Derivation and FEM Simulation}
\label{sec:phase2}

\textbf{Duration:} 8 weeks (June -- July 2026).

\subsubsection{2a: Complete PDE Derivation in \LaTeX}

The full analytical derivation will be typeset with all intermediate steps:
\begin{itemize}[noitemsep]
  \item Maxwell's equations $\to$ vector potential form $\to$ Coulomb gauge reduction
  \item Convective diffusion PDE (\cref{eq:cmd}) in rotor frame
  \item Complex-exponential ansatz, boundary-value problem, and closed-form solution
    for $\hat{A}(y)$
  \item Maxwell Stress Tensor integration yielding \cref{eq:thrust_analytical}
  \item Proof of \cref{prop:opt_slip} with all algebraic steps
  \item Proof of \cref{thm:vol_cancel} (extend current sketch to full detail)
  \item Derivation of the thin-sheet vs.\ skin-effect transition frequency
    $\omega_s^{\mathrm{tr}}$
\end{itemize}

\subsubsection{2b: FEM Implementation and Validation}

\begin{enumerate}[noitemsep, label=\arabic*.]
  \item Implement the 2D transient FEM in FEniCS (Python UFL formulation)
  \item Validate against the analytical solution \cref{eq:thrust_analytical} at
    low Rm (no convection): target $< 5\%$ discrepancy
  \item Execute mesh convergence study (coarse / medium / fine) and document
    Richardson extrapolation bounds
  \item Parametric sweep: $3 \times 5 = 15$ operating points per alloy (30 total)
  \item Post-process: extract $\mathcal{F}^*$ vs.\ $S$ for each operating point and
    overlay on \cref{eq:scaling_law}
  \item Visualise magnetic flux density $\|\mathbf{B}\|$ and eddy current density
    $\mathbf{J}$ at $S = 0.5,\,1.0,\,2.0$ to illustrate skin-effect evolution
\end{enumerate}

\textbf{Go/no-go criterion:} FEM thrust values within $\pm 10\%$ of analytical
prediction at $\mathrm{Rm} \ll 1$. If not achieved, the mesh or boundary conditions
will be revised before proceeding to hardware.

\textbf{Deliverable:} Jupyter notebook (committed to GitHub) with all FEM results,
convergence plots, and parametric sweep figures — equivalent to a journal article's
simulation section.

\subsection{Phase 3 — Hardware Design and Manufacturing}

\textbf{Duration:} 6 weeks (July -- August 2026).

\subsubsection{3a: KiCad PCB Layout}

\begin{enumerate}[noitemsep, label=\arabic*.]
  \item Schematic capture: RP2040, DRV8313, INA240A, TCST2103 array, AS5311, power
    connectors, debug header
  \item Layer 4 (control) layout: component placement observing heat dissipation
    constraints (DRV8313 thermal pad to copper pour)
  \item Layers 1--3 (stator): meandered phase traces for all three pole-pitch
    configurations, with solder-jumper pole selection
  \item Design rule check (DRC): minimum 0.2\,mm clearance (JLCPCB 4-layer standard)
  \item Gerber export and JLCPCB DFM review
\end{enumerate}

\subsubsection{3b: Mechanical Rail Design}

A 3D-printed PETG launch rail will constrain the rotor to 1D translational motion:
\begin{itemize}[noitemsep]
  \item Air gap $g = 1.0 \pm 0.1$\,mm between rotor underside and stator surface
    (controlled by precision shims)
  \item Rail length: 300\,mm (accommodates 8 photointerrupters at 15\,mm pitch
    plus entry/exit clearance)
  \item Rotor guide channels: $\pm 0.2$\,mm lateral clearance to minimise friction
\end{itemize}

\textbf{Deliverable:} KiCad project files and STL files committed to GitHub;
Gerber files submitted to JLCPCB.

\subsection{Phase 4 — Firmware Development and System Integration}

\textbf{Duration:} 5 weeks (August -- September 2026).

\subsubsection{4a: RP2040 Firmware (C++)}

\begin{enumerate}[noitemsep, label=\arabic*.]
  \item PIO SPWM engine: 3-channel, 256-entry sine LUT, configurable frequency
    via DMA period register
  \item Real-time slip control loop: read AS5311 velocity via SPI $\to$ compute
    $\omega_s = \omega - \pi v/\tau_p$ $\to$ update PIO frequency every 1\,ms
  \item DAQ interrupt handler: TCST2103 beam-break timestamps stored in ring buffer
    (8-byte entries: sensor ID + 32-bit timer count at 125\,MHz)
  \item UART telemetry: stream timestamp data to host PC at 921600 baud during
    each experimental run
\end{enumerate}

\subsubsection{4b: Host Software (Python)}

\begin{enumerate}[noitemsep, label=\arabic*.]
  \item Serial parser: reconstruct position--time traces from UART stream
  \item Velocity and acceleration extraction: Savitzky--Golay filter (order 3,
    window 7) applied to position data
  \item Thrust inference: both instantaneous-power and impulse methods implemented
    and compared
  \item Automated data export to \texttt{.csv} with metadata (configuration,
    timestamp, run ID)
\end{enumerate}

\textbf{System integration test:} Before rotor launch, verify three-phase waveform
symmetry with logic analyser, confirm phase current measurement with INA240A vs.\
clamp meter, and verify photointerrupter timing against a reference oscillator.

\textbf{Go/no-go criterion:} Three-phase SPWM THD $< 2\%$; photointerrupter timing
agreement $< 10\,\mu$s between board-side and logic analyser timestamp.

\textbf{Deliverable:} First rotor launch (any motion detected) by end of September 2026.

\subsection{Phase 5 — Parametric Experiment, Scaling Law Fit, and Thesis}

\textbf{Duration:} 10 weeks (October -- December 2026).

\subsubsection{5a: Experimental Data Collection}

Execute the $3 \times 3 \times 2 = 18$-run factorial design (\cref{tab:factorial}).
For each run:
\begin{itemize}[noitemsep]
  \item Set $\tau_p$ via solder jumpers and verify with LCR meter (inductance change)
  \item Set $\omega_s$ via firmware parameter; confirm with logic analyser
  \item Install rotor alloy; measure temperature with thermocouple pre/post run
  \item Acquire 5 repeat launches per operating point (90 total launches)
  \item Record: position--time (both DAQ modalities), phase currents, supply voltage
\end{itemize}

\subsubsection{5b: Scaling Law Fit}

\begin{enumerate}[noitemsep, label=\arabic*.]
  \item Compute $\mathcal{F}^*$ and $S$ for each of the 18 operating points
  \item Plot $\mathcal{F}^*$ vs.\ $S$ with error bars from \cref{tab:uncertainty}
  \item Fit \cref{eq:scaling_law} to data: 1-parameter fit (amplitude $\mathcal{F}^*_{\max}$
    if curve shape is fixed, or 2-parameter if an empirical correction is needed)
  \item Report goodness of fit ($R^2$, residual RMS) and test whether the data
    are consistent with $S^* = 1$ at the $2\sigma$ level
  \item Generate design charts: $\omega_s^*$ vs.\ $\tau_p$ for 7075 and 6061 alloys,
    with FEM and analytical overlays
\end{enumerate}

\subsubsection{5c: Thesis Writing}

The thesis will follow the standard MIT AeroAstro format:
\begin{enumerate}[noitemsep, label=\arabic*.]
  \item Abstract + Introduction (research gap, claim, objectives)
  \item Theoretical Framework (Sections \ref{sec:phase2} content, fully typeset in \LaTeX)
  \item Numerical Simulation (FEM methodology, convergence, parametric sweep)
  \item Experimental Design (PCB architecture, DAQ, uncertainty budget)
  \item Results and Discussion (scaling law fit, comparison with theory and FEM)
  \item Conclusion and Future Work (3D FEM, active slip control, space-launch scaling)
  \item Bibliography, Appendices (full derivations, KiCad files, firmware listing)
\end{enumerate}

\textbf{Final deliverable:} Thesis draft submitted December 2026.

%% ══════════════════════════════════════════════════════════════════════════════
\section{Future Work and Extensions}
%% ══════════════════════════════════════════════════════════════════════════════

The validated scaling law opens several natural extensions, each of which could
constitute a follow-on publication:

\begin{itemize}[noitemsep]
  \item \textbf{3D end effects.} The 2D FEM neglects transverse edge effects at the
    rotor ends. A 3D transient simulation would quantify the end-effect efficiency
    penalty and its dependence on the rotor aspect ratio $L_{\mathrm{rotor}}/\tau_p$.
  \item \textbf{Closed-loop slip control.} The firmware infrastructure supports
    real-time slip tracking. An active controller maintaining $\omega_s = \omega_s^*$
    throughout the acceleration run could be validated against open-loop results,
    quantifying the efficiency gain.
  \item \textbf{Thermal limits.} The INA240A current data enable a thermal model of
    the PCB stator under sustained operation, informing duty-cycle limits for continuous
    (rather than pulsed) LIM applications.
  \item \textbf{Space-launch scaling.} The validated dimensionless scaling law
    \cref{eq:scaling_law} can be extrapolated to lunar mass-driver scales, providing
    a first-principles estimate of the pole pitch and slip frequency required to achieve
    a target specific impulse without an applied magnetic field.
\end{itemize}

%% ══════════════════════════════════════════════════════════════════════════════
\section{Conclusion}
%% ══════════════════════════════════════════════════════════════════════════════

This proposal describes a Masters-level research program at the intersection of
continuum electrodynamics (EE) and electromagnetic launch technology (AeroAstro) that:

\begin{enumerate}[noitemsep]
  \item Makes a \textbf{falsifiable scientific claim}: the dimensionless thrust
    $\mathcal{F}^* = S/(1+S^2)$ is maximised at $S = 1$, equivalently at
    $\omega_s^* = \pi^2/(\mu_0\sigma\tau_p^2)$, and this prediction holds across
    rotor materials and pole-pitch scales.
  \item Provides a \textbf{three-tier validation}: analytical derivation, 2D transient
    FEM with documented mesh convergence, and 18-point factorial experiment with
    formal uncertainty budgets.
  \item Delivers \textbf{aerospace-usable design rules}: charts of $\omega_s^*$ vs.\
    $\tau_p$ for commercially available aluminium alloys, directly applicable to
    sub-100\,W EMALS analog and electromagnetic launch preliminary design.
  \item Is \textbf{executable within \$400} and a 7-month timeline using exclusively
    open-source software (FEniCS, KiCad, Python) and commodity hardware (JLCPCB, DigiKey).
\end{enumerate}

The project is distinguished from prior work not by greater hardware complexity but by
greater intellectual architecture: it transforms a compelling demonstration into a
rigorous scientific result that advances the field's quantitative understanding of
miniature linear electromagnetic accelerators.

%% ══════════════════════════════════════════════════════════════════════════════
%% BIBLIOGRAPHY
%% ══════════════════════════════════════════════════════════════════════════════
\newpage
\begin{thebibliography}{99}

\bibitem{Doyle2010}
M.~R. Doyle, D.~J. Samuel, T.~Conway, and R.~R. Klimowski,
``Electromagnetic Aircraft Launch System --- EMALS,''
\textit{IEEE Trans.\ Magnetics}, vol.~31, no.~1, pp.~528--533, Jan. 1995.

\bibitem{Gieras1994}
J.~F. Gieras,
\textit{Linear Induction Drives}.
Oxford: Clarendon Press, 1994.

\bibitem{Boldea2001}
I.~Boldea and S.~A. Nasar,
\textit{Linear Electric Actuators and Generators}.
Cambridge: Cambridge University Press, 1997.

\bibitem{Davidson2001}
P.~A. Davidson,
\textit{An Introduction to Magnetohydrodynamics}.
Cambridge: Cambridge University Press, 2001.

\bibitem{Griffiths2017}
D.~J. Griffiths,
\textit{Introduction to Electrodynamics}, 4th ed.
Cambridge: Cambridge University Press, 2017.

\bibitem{Kolm1980}
H.~Kolm and R.~Mongeau,
``Basic principles of coaxial launch technology,''
\textit{IEEE Trans.\ Magnetics}, vol.~20, no.~2, pp.~227--230, Mar. 1984.

\bibitem{Lewin2006}
P.~L. Lewin, J.~E. Theed, A.~E. Davies, and S.~T. Larsen,
``Method for rating power transformers with non-sinusoidal currents,''
\textit{IEE Proc.\ Electr.\ Power Appl.}, vol.~146, no.~4, pp.~411--416, 1999.

\bibitem{Raminosoa2018}
T.~Raminosoa, J.~A. Farooq, A.~Djerdir, and A.~Miraoui,
``Reluctance network modelling of surface permanent magnet motor considering iron
nonlinearities,''
\textit{Energy Conversion and Management}, vol.~50, no.~5, pp.~1356--1361, 2009.

\bibitem{Maloberti2020}
O.~Maloberti, V.~Mazauric, G.~Meunier, A.~Kedous-Lebouc, O.~Geoffroy, Y.~Rebiere,
P.~Labie, and Y.~Marchand,
``An energy-based formulation for dynamic hysteresis and extra-losses,''
\textit{IEEE Trans.\ Magnetics}, vol.~42, no.~4, pp.~895--898, Apr. 2006.

\bibitem{Connor2015}
J.~Connor, M.~Dommer, E.~Merritt, and C.~Grantham,
``EMALS system integration and testing,''
in \textit{Proc.\ IEEE Electric Ship Technologies Symposium (ESTS)},
Alexandria, VA, Jun. 2015, pp.~28--34.

\bibitem{AlnaesEtal2015}
M.~S.~Aln\ae{}s et al.,
``The FEniCS Project Version 1.5,''
\textit{Archive of Numerical Software}, vol.~3, no.~100, pp.~9--23, 2015.

\end{thebibliography}

\end{document}